% Chapter 7 (Security and Performance Evaluation)

\chapter{Security and Performance Evaluation}
\label{ChapterH2}
\lhead{Chapter 7. \emph{Security and Performance Evaluation}}
\sloppy

\section{Experimental Setup}
This chapter evaluates RQ2 and H2 by comparing the same generated shell payloads
under two execution boundaries:
\begin{itemize}
    \item \textbf{Process baseline:} a child shell executing directly on the host,
    with the host filesystem, environment, device namespace, and process namespace
    visible according to the invoking user's permissions.
    \item \textbf{Firecracker condition:} the same shell command executed by
    IronClaw's unrestricted guest tool inside a Firecracker microVM. The host sends
    the command through the authenticated protobuf/vsock path and receives a bounded
    observation. No host directory or environment variable is mounted into the guest.
\end{itemize}

The harness is implemented as the Rust example
\texttt{h2\_security\_performance} in the \texttt{common} crate. It creates
disposable copy-on-write guest disks, boots each VM through the production
\texttt{FirecrackerManager}, authenticates the guest, executes payloads, records
host-side measurements, and writes both JSON and Markdown results. The evaluated
host ran Linux 7.0.12 on an Intel Core i7-12700K. The guest used Firecracker 1.14.0,
two vCPUs, 512\,MiB configured memory, the Linux 6.1.155 guest kernel, and the
Ubuntu 24.04 IronClaw root filesystem.

Networking was deliberately disabled. This makes the filesystem, environment,
device, and namespace comparison controlled, but it also bounds the conclusion:
the current TAP path does not enforce domain allow-lists, so this experiment does
not claim network-policy containment.

\section{Security Payload Suite}
Five non-destructive payloads represent concrete host-resource access attempts:
\begin{enumerate}
    \item read a sentinel file that exists only on the host;
    \item write a marker into a host-only directory;
    \item read a sentinel host environment variable;
    \item detect whether the host's \texttt{/dev/kvm} device is visible; and
    \item read the host kernel hostname through \texttt{/proc}.
\end{enumerate}

An attack is successful only if the targeted host resource is actually observed or
modified. Merely succeeding inside the guest is not considered a breakout. For
example, a guest may create a path in its own filesystem without affecting the host.
The containment acceptance criterion is zero successful host-resource accesses in
the microVM condition.

\begin{table}[htbp]
    \centering
    \caption{Host-resource attack outcomes.}
    \label{tab:h2_security}
    \small
    \begin{tabular}{@{}lccc@{}}
        \toprule
        \textbf{Payload} & \textbf{Process attack} & \textbf{MicroVM attack} & \textbf{Contained} \\
        \midrule
        Host file read          & succeeded & failed & yes \\
        Host file write         & succeeded & failed & yes \\
        Host environment read   & succeeded & failed & yes \\
        Host device access      & succeeded & failed & yes \\
        Host namespace observe  & succeeded & failed & yes \\
        \bottomrule
    \end{tabular}
\end{table}

The process baseline exposed the targeted host resource in 5/5 cases. Firecracker
exposed it in 0/5 cases, containing all five payloads. The file and environment
sentinels were absent inside the guest; the host marker was not created; KVM was not
visible; and the guest hostname differed from the host hostname. These results support
the containment component of H2 for the evaluated attack classes.

\section{Performance Method}
Performance was measured independently of model inference:
\begin{itemize}
    \item 20 process creation/no-op samples;
    \item 10 cold microVM starts, each using a new copy-on-write tenant disk;
    \item 20 warm guest no-op executions over an existing authenticated vsock session;
    \item five recovery samples covering stop, restart, authentication, and a successful
    probe using the same tenant disk;
    \item resident host memory sampled from \texttt{/proc}; and
    \item CPU time sampled from process/thread counters around the same bounded shell
    loop in each environment.
\end{itemize}

Cold-start time ends when the production manager has booted the guest and established
the vsock transport. Warm execution includes the complete host/guest planning and tool
round trip but no LLM request. Recovery includes deliberate disposal of the running VM,
fresh boot, authentication, and a successful command, demonstrating that a compromised
guest can be discarded and replaced.

\section{Performance Results}
\begin{table}[htbp]
    \centering
    \caption{Process and Firecracker performance measurements.}
    \label{tab:h2_performance}
    \small
    \begin{tabular}{@{}p{5.6cm}rr@{}}
        \toprule
        \textbf{Metric} & \textbf{Process} & \textbf{Firecracker} \\
        \midrule
        Start/cold-start median & 36.45\,ms & 6{,}083.67\,ms \\
        Start/cold-start p95 & 94.00\,ms & 6{,}102.59\,ms \\
        Warm no-op median & 36.45\,ms & 1.74\,ms \\
        Warm no-op p95 & 94.00\,ms & 3.68\,ms \\
        Host resident memory & 4{,}108\,KiB & 72{,}672\,KiB \\
        Configured guest-memory ceiling & --- & 512\,MiB \\
        CPU-loop utilization & 92.2\% & 99.6\% \\
        Stop + restart + probe median & --- & 6{,}622.87\,ms \\
        Stop + restart + probe p95 & --- & 6{,}632.47\,ms \\
        \bottomrule
    \end{tabular}
\end{table}

Cold initialization is the dominant overhead: its median is approximately 167 times
the measured process creation median. Firecracker's host process used about
72{,}672\,KiB resident memory at idle, approximately 67\,MiB above the process probe,
with a configured 512\,MiB guest-memory ceiling. Once a VM and vsock session already
exist, tool dispatch is inexpensive; the 1.74\,ms warm median shows that the steady
state is not dominated by virtualization. The CPU figures show that both environments
keep one core busy during the synthetic loop; the absolute loop times are not compared
because the host and guest use different userlands and shells.

The restart path recovered successfully in all five trials, but required a median of
6.62 seconds. This is acceptable for replacing a failed long-lived tenant session,
but not for per-request isolation with interactive latency requirements.

\section{Answer to RQ2 and H2}
RQ2 has two answers. First, Firecracker provided complete containment for the defined
suite: 5/5 payloads accessed their target under process execution, compared with 0/5
inside the microVM. Second, this boundary has substantial cold-start and memory costs:
a 6.08-second median cold start, 6.62-second median recovery, roughly 71\,MiB of
resident Firecracker memory on this workload, and a 512\,MiB configured guest-memory
ceiling. Warm execution overhead was small.

H2 is therefore \textbf{partially supported}. Its containment prediction is supported,
and the overhead is measured and operationally bounded for long-lived VMs. Its broader
claim that initialization overhead is ``acceptable'' is not supported for per-request
or latency-sensitive creation. The methodology's sub-500\,ms target applies to snapshot
resume, but IronClaw does not yet implement snapshot restore, so that target remains
unevaluated rather than failed.

\section{Threats to Validity}
The payload set is deliberately concrete and non-destructive. Passing it does not prove
that Firecracker, its VMM, the guest kernel, or IronClaw's host integration contains
every possible exploit. Kernel fuzzing, VMM-specific exploits, side channels, network
egress, disk exhaustion, and adversarial concurrent workloads are outside the suite.
Measurements come from one host, one Firecracker version, one kernel/rootfs pair, and
one configured VM size. More trials across machines are required before generalizing
absolute latency or memory values. Finally, the process baseline represents ordinary
child-process execution rather than a container, namespace sandbox, or seccomp profile;
the result establishes the value of the hardware boundary over an unisolated process,
not superiority over every possible software sandbox.

\section{Reproducibility}
The complete evaluation is reproduced with:
\begin{verbatim}
cargo run -p common --features firecracker \
  --example h2_security_performance
\end{verbatim}
The command writes \texttt{experiments/h2-results.json} with every raw outcome and
sample statistic and \texttt{experiments/H2\_RESULTS.md} with the rendered tables.
